\section{Simulações Numéricas e Análise dos Resultados}
Para validar a transferência de Hohmann, as equações do movimento absoluto, do movimento relativo e a lei de controle proposta, foram realizadas simulações numéricas em ambiente MATLAB/Simulink.
O objetivo desta etapa é avaliar a resposta dinâmica do sistema frente a condições iniciais estabelecidas e a diferentes níveis de complexidade na modelagem.

\subsection{Configuração inicial}
As simulações consideram a definição das condições iniciais de atitude tanto do \emph{chaser} quanto do \emph{target}. O primeiro é modelado com desvios angulares e velocidades iniciais que representam cenários realistas de aproximação, enquanto o segundo é tratado como referência, podendo estar em repouso no referencial inercial ou com orientação fixa no referencial local LVLH.

\subsection{Integração numérica}
A evolução temporal das equações diferenciais de movimento é obtida por integração numérica. Para isso, são utilizados algoritmos de passo variável disponíveis no MATLAB/Simulink, que garantem a precisão desejada na propagação dos estados e estabilidade no tratamento de equações não lineares. Essa abordagem permite avaliar a robustez do controle PD em diferentes condições de simulação.

\subsection{Casos de teste}
Diversos casos de teste são concebidos para analisar o desempenho do sistema como um todo, sendo:

(i) transferência orbital do \emph{chaser} para a órbita desejada do \emph{target} utilizando a transferência de Hohmann;

(ii) controle absoluto da atitude no referencial inercial, verificando a capacidade do controlador em estabilizar o \emph{chaser} independentemente da referência externa;

(iii) controle relativo no cenário entre \emph{chaser} e \emph{target}, em que a atitude do \emph{chaser} deve acompanhar a orientação do \emph{target} durante a aproximação;

(iv) comparação de erros residuais, com análise do desvio entre as atitudes simuladas e as atitudes de referência.

\subsection{Avaliação de desempenho}
A análise dos resultados se baseia nos seguintes critérios:

(i) estabilização, isto é, a convergência da atitude para a orientação desejada sem oscilações exageradas ou não físicas;

(ii) robustez, medida pela capacidade de manter desempenho aceitável frente a variações inerciais provocadas pelo movimento do manipulador e por perturbações externas;

(iii) saturação de atuadores, que corresponde à verificação se os torques exigidos pelo controlador permanecem dentro dos limites físicos do sistema de acionamento;

(iv) tempo de resposta, avaliado pela duração necessária para que o erro angular relativo esteja de acordo com os critérios estabelecidos;

(v) atendimento às condições de sincronização posição--atitude, em especial erro angular inferior a $1^\circ$ e permanência do \emph{chaser} dentro do corredor seguro de aproximação, conforme estabelecido na seção \eqref{sec:sincronizacao}.

Assim, as simulações numéricas fornecem a ligação entre a modelagem matemática do sistema e sua análise de desempenho, permitindo verificar a validade da estratégia de controle adotada e a aderência aos critérios de sincronização posição--atitude.
